\documentclass[11pt,a4paper]{article}
\usepackage[utf8]{inputenc}
\usepackage[T1]{fontenc}
% \usepackage[polish]{babel}  % wymaga texlive-lang-polish
\usepackage[margin=2.5cm]{geometry}
\usepackage{lmodern}
\usepackage{parskip}
\usepackage{microtype}
\usepackage{booktabs}
\usepackage{longtable}
\usepackage{listings}
\usepackage{xcolor}
\usepackage{hyperref}
\usepackage{enumitem}
\usepackage{titlesec}

\hypersetup{
    colorlinks=true,
    linkcolor=blue!60!black,
    urlcolor=blue!60!black,
    citecolor=blue!60!black,
}

\definecolor{codebg}{RGB}{245,245,245}
\definecolor{codecomment}{RGB}{100,100,100}
\definecolor{codestring}{RGB}{0,130,0}
\definecolor{codekeyword}{RGB}{0,0,180}

\lstdefinestyle{bash}{
    backgroundcolor=\color{codebg},
    basicstyle=\ttfamily\small,
    commentstyle=\color{codecomment},
    keywordstyle=\color{codekeyword},
    stringstyle=\color{codestring},
    breaklines=true,
    frame=single,
    framerule=0.4pt,
    rulecolor=\color{gray!50},
    xleftmargin=6pt,
    xrightmargin=6pt,
    aboveskip=6pt,
    belowskip=6pt,
}
\lstset{style=bash}

\titleformat{\section}{\large\bfseries}{}{0pt}{}[\titlerule]
\titleformat{\subsection}{\normalsize\bfseries}{}{0pt}{}

\title{\textbf{io-2026s --- Instrukcja użytkownika}\\[4pt]
    \large Harmonogramowanie na jednej maszynie z optymalizacją hiperparametrów przez LLM}
\author{Karolina Piotrowska, Jan Rosa}
\date{\today}

\begin{document}
\maketitle
\tableofcontents
\newpage

% ---------------------------------------------------------------------------
\section{Przegląd}
% ---------------------------------------------------------------------------

\textbf{io-2026s} stosuje duże modele językowe (LLM) oraz klasyczne metaheurystyki
do rozwiązywania \emph{problemu harmonogramowania na jednej maszynie ze wspólnym
terminem realizacji} (instancje ORLib \texttt{sch*}).
Agenty przeszukują przestrzeń permutacji zadań, minimalizując łączny ważony
koszt zbyt wczesnego i zbyt późnego zakończenia:

\[
  C(\pi) = \sum_i \bigl[ a_i \cdot \max(0,\, d - C_i) + b_i \cdot \max(0,\, C_i - d) \bigr]
\]

gdzie $d = \lfloor h \cdot \sum p_i \rfloor$ jest wspólnym terminem realizacji
sterowanym współczynnikiem szczelności $h$, a $C_i$ to czas ukończenia zadania~$i$.

% ---------------------------------------------------------------------------
\section{Instalacja}
% ---------------------------------------------------------------------------

\subsection{Wymagania}

\begin{itemize}[noitemsep]
  \item Python 3.10+
  \item Menedżer pakietów \texttt{uv} --- \url{https://docs.astral.sh/uv/}
  \item GCC (do kompilacji rozszerzeń C)
  \item Ollama (tylko dla agentów GEPA) --- \url{https://ollama.com/}
\end{itemize}

\subsection{Konfiguracja środowiska}

\begin{lstlisting}[language=bash]
# Sklonuj repozytorium, zainstaluj zaleznosci i skompiluj rozszerzenia C
uv sync
\end{lstlisting}

Polecenie \texttt{uv sync} instaluje pakiety Pythona, buduje rozszerzenia C
(\texttt{evaluate\_opt.c}, \texttt{crossover\_opt.c}) i zapisuje skompilowane
pliki \texttt{.so} w katalogu
\texttt{src/native\_optimized/\_\_compiled\_\_/}.

% ---------------------------------------------------------------------------
\section{Główny skrypt benchmarkowy --- \texttt{main.py}}
% ---------------------------------------------------------------------------

\subsection{Podstawowe wywołanie}

\begin{lstlisting}[language=bash]
python main.py [OPCJE]
\end{lstlisting}

Uruchamia wybrane agenty na zbiorze danych, drukuje tabelę wyników ASCII,
zapisuje \texttt{benchmark.csv} i (o ile nie podano \texttt{--no-plots})
generuje wykresy zbieżności oraz porównawcze.
Wszystkie wyniki trafiają do katalogu \texttt{results/} w podkatalogu
oznaczonym znacznikiem czasowym.

\subsection{Parametry}

\begin{longtable}{@{}p{4.5cm} p{2.2cm} p{2.2cm} p{6.3cm}@{}}
\toprule
\textbf{Flaga} & \textbf{Typ} & \textbf{Domyślnie} & \textbf{Opis} \\
\midrule
\endhead

\texttt{--data} & ścieżka & \texttt{data/sch10.txt} &
    Plik zbioru danych ORLib. Dostępne pliki: \texttt{sch10}, \texttt{sch20},
    \texttt{sch50}, \texttt{sch100}, \texttt{sch200}, \texttt{sch500},
    \texttt{sch1000}, \texttt{schpascal}. \\[4pt]

\texttt{--n-instances} & int & 5 &
    Liczba instancji problemu do zbenchmarkowania
    (pobierane od początku zbioru danych). \\[4pt]

\texttt{--h} & float & 0.4 &
    Współczynnik szczelności terminu realizacji. Mniejsze wartości
    utrudniają problem (większość zadań spóźniona), większe go upraszczają.
    Typowe wartości ORLib: 0.2, 0.4, 0.6, 0.8. \\[4pt]

\texttt{--max-steps} & int & 50 &
    Długość epizodu (liczba ruchów zamiany) dla agentów iteracyjnych.
    Podczas treningu GEPA stosowana jest domyślna wartość $10 \times n$,
    gdy parametr nie jest ustawiony. \\[4pt]

\texttt{--seed} & int & 42 &
    Ziarno generatora liczb losowych (dla powtarzalności wyników). \\[4pt]

\texttt{--out-dir} & ścieżka & \texttt{results} &
    Katalog wyjściowy. Automatycznie tworzony jest podkatalog ze znacznikiem,
    np.\ \texttt{results/sch10\_sa-genetic\_h0.4\_n5\_s50\_20260607\_120000/}. \\[4pt]

\texttt{--no-plots} & flaga & wył. &
    Pomija generowanie wykresów matplotlib.
    Przydatne w środowiskach bez wyświetlacza lub w CI. \\[4pt]

\texttt{--agents} & string & \texttt{random,greedy,sa,genetic} &
    Lista agentów rozdzielona przecinkami. Dostępne klucze: \texttt{random},
    \texttt{greedy}, \texttt{sa}, \texttt{genetic}, \texttt{gepa-sa},
    \texttt{gepa-ga}. \\[4pt]

\texttt{--train-instances} & int & 3 &
    \textit{(tylko GEPA)} Liczba instancji używanych podczas przeszukiwania
    hiperparametrów przez GEPA. Odrębny zbiór od zestawu benchmarkowego. \\[4pt]

\texttt{--max-metric-calls} & int & 50 &
    \textit{(tylko GEPA)} Łączny budżet ewaluacji optymalizatora GEPA
    (zlicza każde uruchomienie agent$\times$instancja). \\[4pt]

\texttt{--reflection-lm} & string & \texttt{ollama/qwen3:4b} &
    \textit{(tylko GEPA)} Identyfikator modelu LiteLLM używanego przez LLM
    do refleksji. Można użyć dowolnego modelu Ollama,
    np.\ \texttt{ollama/gemma4:e2b}. \\[4pt]

\texttt{--interactions-log} & ścieżka & \textit{brak} &
    \textit{(tylko GEPA)} Jeśli podano, wszystkie surowe pary prompt/odpowiedź
    LLM są zapisywane do tego pliku JSON po zakończeniu treningu.
    Sama nazwa pliku (bez katalogu) umieszcza go w oznaczonym podkatalogu
    wyjściowym. \\[4pt]

\texttt{-v / --verbose} & flaga & wył. &
    Wyświetla postęp dla każdej instancji podczas benchmarkowania. \\

\bottomrule
\end{longtable}

\subsection{Przykładowe wywołania}

\begin{lstlisting}[language=bash]
# Szybki test: wszystkie klasyczne agenty na instancjach 10-zadaniowych
python main.py --agents random,greedy,sa,genetic --data data/sch10.txt

# Pelny benchmark z GEPA, bez wykresow
python main.py \
    --agents gepa-ga,genetic,gepa-sa,sa \
    --data data/sch1000.txt \
    --n-instances 5 \
    --train-instances 5 \
    --max-metric-calls 100 \
    --reflection-lm ollama/qwen3:4b-instruct-2507-q4_K_M \
    --interactions-log interactions.json \
    --no-plots

# Dokladna reprodukcja eksperymentu
python main.py --seed 0 --h 0.6 --max-steps 200
\end{lstlisting}

\subsection{Struktura katalogu wynikowego}

\begin{lstlisting}
results/<znacznik>/
    benchmark.csv              # srednia/odch. std poprawy% dla kazdego agenta
    gepa_history_GEPA_GA.png   # wykres wyniku GEPA w kolejnych iteracjach
    gepa_history_GEPA_SA.png
    convergence.png            # koszt w funkcji kroku dla kazdego agenta
    agent_comparison.png       # wykres slupkowy wynikow koncowych
    cost_breakdown.png         # wczesnosc vs. spoznienie
    schedule_comparison.png    # harmonogram przed i po optymalizacji (Gantt)
    interactions.json          # surowe prompty LLM (jesli zadan)
\end{lstlisting}

% ---------------------------------------------------------------------------
\section{Dostępne agenty}
% ---------------------------------------------------------------------------

\begin{longtable}{@{}p{3.8cm} p{2.5cm} p{8.2cm}@{}}
\toprule
\textbf{Agent} & \textbf{Klucz CLI} & \textbf{Opis} \\
\midrule
\endhead
RandomAgent        & \texttt{random}  & Losowy wybór zamiany na każdym kroku. \\[4pt]
GreedyAgent        & \texttt{greedy}  & Wyczerpujące przeszukiwanie 1-krokowe;
                                        wybiera najlepszą pojedynczą zamianę. \\[4pt]
SimulatedAnnealing & \texttt{sa}      & Wyżarzanie symulowane z oscylującym
                                        harmonogramem temperatury:
                                        $T(t) = T_0 e^{-bt} + c\sin^2(t/d)$.
                                        Hiperparametry: $T_0$, $b$, $c$, $d$, $T_{\min}$. \\[4pt]
GeneticAlgorithm   & \texttt{genetic} & Krzyżowanie OX z elityzmem.
                                        Ponowna optymalizacja co
                                        \texttt{reoptimize\_every} kroków.
                                        Hiperparametry: rozmiar populacji,
                                        liczba pokoleń, współczynnik mutacji,
                                        rozmiar elity. \\[4pt]
GEPA-SA            & \texttt{gepa-sa} & Używa GEPA (przeszukiwanie stylem
                                        bayesowskim sterowane przez LLM) do
                                        znalezienia optymalnych hiperparametrów SA,
                                        a następnie uruchamia SA z tymi parametrami. \\[4pt]
GEPA-GA            & \texttt{gepa-ga} & Analogicznie jak wyżej, lecz dla
                                        hiperparametrów GA. \\
\bottomrule
\end{longtable}

\subsection{Działanie agenta GEPA}

GEPA ocenia kandydackie konfiguracje hiperparametrów, uruchamiając bazowego
agenta na podzbiorze treningowym i raportując procentową poprawę.
Po każdej ewaluacji LLM proponuje udoskonaloną konfigurację.
Konsola drukuje jeden blok na iterację GEPA:

\begin{lstlisting}
=== Iter 1  ACCEPTED  [*** NEW BEST ***]  ========================
  Config  : population_size=50  generations=5  mutation_rate=0.1 ...
> SCORE  : 0.415  (D+0.096 vs curr-prog)  [GEPA optimises this]
  Quality : improv%  [61.9%  71.4%  63.0%  36.6%  61.1%]  mean=58.8%
  Valset  : composite avg 0.393  [0.422  0.522  0.434  0.166  0.421]
  Time    : 0.09 s/run   Calls: 16 / 100
==================================================================
\end{lstlisting}

\textbf{SCORE} (złożony = jakość $-$ kara za czas) jest miarą optymalizowaną
przez GEPA; \textbf{Quality} pokazuje surowe procenty poprawy dla każdej instancji.

% ---------------------------------------------------------------------------
\section{Przeszukiwanie wielu modeli --- \texttt{run\_sweep.py}}
% ---------------------------------------------------------------------------

Uruchamia \texttt{main.py} sekwencyjnie dla każdego modelu z zakodowanej listy,
najpierw sprawdza dostępność w Ollama, przekierowuje wyjście do pliku logu
i drukuje tabelę podsumowującą.

\begin{lstlisting}[language=bash]
uv run python run_sweep.py
\end{lstlisting}

Edytuj górną część pliku \texttt{run\_sweep.py}, aby skonfigurować przeszukiwanie:

\begin{lstlisting}[language=Python]
MODELS = [
    "ollama/qwen3:4b-instruct-2507-q4_K_M",
    "ollama/gemma4:e2b",
    "ollama/phi4-mini-reasoning:latest",
    "ollama/granite4:1b-h-q8_0",
]

BASE_ARGS = {
    "--agents":           "gepa-ga,genetic,gepa-sa,sa",
    "--data":             "data/sch1000.txt",
    "--train-instances":  "5",
    "--max-metric-calls": "100",
    "--interactions-log": "interactions.json",
}
\end{lstlisting}

Pliki logów są zapisywane do
\texttt{results/sweep\_<model>\_<znacznik>.log}.
Przeszukiwanie kontynuuje działanie nawet w przypadku błędu modelu;
podsumowanie końcowe pokazuje \texttt{OK}, \texttt{FAILED} lub
\texttt{SKIPPED (not installed)} dla każdego modelu.

% ---------------------------------------------------------------------------
\section{Inspektor interakcji z LLM --- \texttt{render\_interactions.py}}
% ---------------------------------------------------------------------------

Renderuje surowy JSON wyprodukowany przez \texttt{--interactions-log}
do czytelnego zestawienia promptów i odpowiedzi.

\begin{lstlisting}[language=bash]
# Wyswietl w terminalu
python render_interactions.py results/<znacznik>/interactions.json

# Zapisz do pliku tekstowego
python render_interactions.py results/<znacznik>/interactions.json \
    --out results/<znacznik>/interactions.txt
\end{lstlisting}

\begin{longtable}{@{}p{4.0cm} p{10.0cm}@{}}
\toprule
\textbf{Argument} & \textbf{Opis} \\
\midrule
\endhead
\texttt{input} (pozycyjny) & Ścieżka do pliku \texttt{interactions.json}. \\[4pt]
\texttt{--out SCIEZKA}     & Zapisuje wyrenderowany tekst do podanej ścieżki
                             zamiast na standardowe wyjście. \\
\bottomrule
\end{longtable}

Każdy blok interakcji zawiera znacznik czasu, nazwę modelu, opóźnienie,
pełny prompt (z szablonem refleksji GEPA, bieżącą konfiguracją i wynikami
ewaluacji) oraz surową odpowiedź JSON modelu.

% ---------------------------------------------------------------------------
\section{Generator dokumentacji --- \texttt{build\_docs.sh}}
% ---------------------------------------------------------------------------

Buduje referencję API w formacie HTML z wyszukiwarką oraz tę instrukcję
użytkownika w formacie PDF.

\begin{lstlisting}[language=bash]
./build_docs.sh
\end{lstlisting}

Wyniki:
\begin{itemize}[noitemsep]
  \item \texttt{docs/\_build/html/index.html} --- HTML Sphinxa z pełnotekstowym
        wyszukiwaniem (motyw Furo).
  \item \texttt{docs/\_build/io-2026s-api.pdf} --- PDF referencji API przez
        LaTeX.
  \item \texttt{docs/\_build/io-2026s-userguide.pdf} --- Niniejsza instrukcja
        użytkownika.
\end{itemize}

Skrypt uruchamia trzy kompilacje Sphinxa (HTML, jednostronicowy HTML, LaTeX)
oraz osobną kompilację \texttt{pdflatex} dla tej instrukcji.

% ---------------------------------------------------------------------------
\section{Generator diagramów UML --- \texttt{generate\_uml.py}}
% ---------------------------------------------------------------------------

Generuje diagram klas UML z pakietu \texttt{src/} przy użyciu biblioteki
\texttt{graphviz}.

\begin{lstlisting}[language=bash]
uv run python generate_uml.py
\end{lstlisting}

Parsuje wszystkie pliki \texttt{.py} w \texttt{src/} za pomocą modułu
\texttt{ast}, wyodrębnia hierarchie klas i sygnatury metod, a następnie
zapisuje \texttt{uml\_diagram.png} (oraz źródło \texttt{.gv}) w katalogu
głównym projektu.

% ---------------------------------------------------------------------------
\section{Profilowanie}
% ---------------------------------------------------------------------------

\begin{lstlisting}[language=bash]
# Szczegolowy profil wywolan -> wykres plomienia HTML
uv run pyinstrument -r html -o logs.html \
    main.py --agents sa,genetic --data data/sch1000.txt --max-steps 10000

# Profilowanie probkujace (niski narzut) -> wykres plomienia SVG
uv run py-spy record -r 500 -o profile.svg -- \
    python main.py --agents genetic --data data/sch1000.txt --max-steps 20000
\end{lstlisting}

Używaj \texttt{pyinstrument} do szczegółowej analizy wywołań;
\texttt{py-spy} do profilowania czasu rzeczywistego bez modyfikowania kodu źródłowego.

% ---------------------------------------------------------------------------
\section{Testy}
% ---------------------------------------------------------------------------

\begin{lstlisting}[language=bash]
# Uruchom wszystkie testy (ok. 2 s)
python -m pytest sanity/ -v

# Jeden plik
python -m pytest sanity/sanity_classical_agents.py -v

# Jeden test
python -m pytest sanity/sanity_agent.py::test_greedy_agent_solve -v
\end{lstlisting}

Testy korzystają z \texttt{data/sch10.txt} (10 zadań) z małą liczbą kroków
(20--30) dla zachowania szybkości. Cztery zestawy: poprawność środowiska,
interfejs agenta, runner benchmarkowy oraz poprawność SA/GA.

% ---------------------------------------------------------------------------
\section{Jakość kodu}
% ---------------------------------------------------------------------------

\begin{lstlisting}[language=bash]
# Lintowanie (z automatycznym naprawianiem)
uv run ruff check src/ sanity/

# Formatowanie
uv run ruff format src/ sanity/
\end{lstlisting}

Ruff jest skonfigurowany w \texttt{pyproject.toml} z wymuszaniem docstringów
w stylu NumPy (reguły \texttt{D}/\texttt{DOC}), bug-bear, pylint oraz
sortowaniem importów.

% ---------------------------------------------------------------------------
\section{Częste problemy}
% ---------------------------------------------------------------------------

\begin{description}[style=nextline]
  \item[Rozszerzenia C nie kompilują się]
    Upewnij się, że GCC jest zainstalowany (\texttt{gcc --version}).
    Na Ubuntu: \texttt{sudo apt-get install build-essential}.
    Następnie ponownie uruchom \texttt{uv sync}.

  \item[GEPA zawiesza się / brak odpowiedzi Ollama]
    Uruchom \texttt{ollama list}, aby sprawdzić, czy model jest pobrany.
    \texttt{run\_sweep.py} sprawdza dostępność automatycznie.

  \item[Wyjście podprocesu zawiesza się w \texttt{run\_sweep.py}]
    \texttt{main.py} jest wywoływany z \texttt{python -u} (niebuforowane),
    co zapewnia strumieniowanie linii po linii --- jest to już zawarte w skrypcie.

  \item[Testy nie mogą znaleźć importów]
    Aktywuj środowisko wirtualne:
    \texttt{source .venv/bin/activate} przed uruchomieniem \texttt{pytest}.
\end{description}

\end{document}
